\label{der:3-2}
设行阶梯形为$R=EA$，$E$可逆。对任何系数列$c$，$Rc=0\Longleftrightarrow Ac=0$，因此每一个列关系的系数在变换前后完全相同。阶梯形的主元列无关：从最后一个主元行向上考察，关系中的最后一个主元列系数先为零，再依次使所有系数为零。主元数为$r$，非主元列可通过从下向上的主元方程唯一表示为这$r$列的组合，因此它们生成全部列。把这些关系返回$A$，得到原主元列组成极大无关组。

若$\mathcal A$可由$\mathcal B$表示，则$U=\Span\mathcal A\subseteq W=\Span\mathcal B$。取$U$的基，它在$W$中仍无关，故$\dim U\le\dim W$。若维数相等而包含为真，取$w\in W\setminus U$，它能加到$U$的基中形成更长无关组，与同维矛盾；所以$U=W$，两组互相可表示。若只知维数相等却没有包含，两条不同直线就是反例。
